\documentclass{article}

% Language setting
\usepackage[english]{babel}

% Set page size and margins
\usepackage[letterpaper,top=2cm,bottom=2cm,left=3cm,right=3cm,marginparwidth=1.75cm]{geometry}

% Useful packages
\usepackage{amsmath}
\usepackage{amssymb}
\usepackage{graphicx}
\usepackage[colorlinks=true, allcolors=black]{hyperref}
\usepackage{titlesec}
\usepackage{xparse}

\NewDocumentCommand{\qcq}{m o o m}{%
	\ensuremath{\forall #1 \IfValueT{#2}{, #2} \IfValueT{#3}{, #3}  \in  \mathbb{#4}}%
}
\NewDocumentCommand{\ext}{m o o m}{%
	\ensuremath{\exists #1 \IfValueT{#2}{, #2} \IfValueT{#3}{, #3}  \in  \mathbb{#4}}%
}

\title{\textbf{Exercices : Probabilités et Dénombrement}}
\author{\textbf{Yassin Akkab}}

\newcommand{\R}{\mathbb{R}}
\newcommand{\N}{\mathbb{N}}
\newcommand{\Z}{\mathbb{Z}}
\newcommand{\C}{\mathbb{C}}
\newcommand{\Q}{\mathbb{Q}}
\newcommand{\D}{\mathbb{D}}

\begin{document}
	
	\maketitle
	
	\subsection*{Exercice 1 (Dénombrement)}
	Une urne contient 5 boules rouges, 3 boules vertes et 2 boules bleues. On tire simultanément 3 boules de l'urne.
	\begin{enumerate}
		\item Quel est le nombre total de tirages possibles ?
		\item Combien de tirages contiennent 3 boules de la même couleur ?
		\item Combien de tirages contiennent au moins une boule rouge ?
	\end{enumerate}
	
	\subsection*{Exercice 2 (Probabilités conditionnelles)}
	Dans une population, 40\% des individus sont vaccinés contre une maladie. La probabilité qu'un individu vacciné tombe malade est de 0,05. La probabilité qu'un individu non vacciné tombe malade est de 0,3.
	On choisit un individu au hasard.
	\begin{enumerate}
		\item Quelle est la probabilité que cet individu soit malade ?
		\item Sachant que l'individu choisi est malade, quelle est la probabilité qu'il soit vacciné ?
	\end{enumerate}
	
	\subsection*{Exercice 3 (Variable aléatoire et loi binomiale)}
	On lance 5 fois de suite un dé équilibré à 6 faces. On s'intéresse au nombre de fois où le dé affiche la face 6.
	Soit $X$ la variable aléatoire représentant ce nombre.
	\begin{enumerate}
		\item Justifier que la variable aléatoire $X$ suit une loi binomiale et préciser ses paramètres.
		\item Calculer la probabilité d'obtenir exactement deux fois la face 6.
		\item Calculer l'espérance mathématique de $X$.
	\end{enumerate}
	
\end{document}
